\section{Lesson 7}

\subsection{用 Bergman 度量定义映射}
回顾上次课, 给定一个黎曼曲面 $\Sigma=(S,J)$, 我们构造了一个具体的微分同胚
\begin{equation*}
    \Phi: \mathrm{Teich}(S)\cong H^0(\Sigma,K^2).
\end{equation*}
借由这个微分同胚, 我们构造了 Teichm\"uller 空间的紧化.
\begin{remark}
    这个紧化不依赖于基点 $\Sigma$ 的选择.
\end{remark}
\begin{remark}
    这个具体微分同胚的构造需要借由从 $\Sigma$ 到 $(S,\rho)$ 的调和映照, 其中 $\rho$ 是
    $S$ 上取定的一个双曲度量. 注意定义调和映照需要目标空间的度量, 现在如果我们不考虑双曲度量, 
    而是与复结构相关的度量, 那么这套过程能走下去吗?
\end{remark}
我们现在考虑黎曼曲面 $\Sigma=(S,J)$ 上的 Bergman 度量, 它的定义如下:

......



\subsection{秩 3 Higgs 丛, Hitchin 分支}
结束上面的讨论, 让我们转到一个新的话题: $n=3$ 时的 Higgs 丛. 
现在我们考虑曲面基本群的表示
\begin{equation*}
    \rho: \pi_1(S)\to \mathrm{PSL}(3,\mathbb{R}).
\end{equation*}
\begin{remark}
    注意 $\mathrm{PSL}(3,\mathbb{R})\cong \mathrm{SL}(3,\mathbb{R})$. 因为 $-I\notin \mathrm{SL}(3,\mathbb{R})$. 
\end{remark}
在 $\mathrm{Hom}(\pi_1(S), \mathrm{PSL}(3,\mathbb{R}))$ 空间中有一类特殊的表示, 它们本质上来自于 $\mathrm{PSL}(2,\mathbb{R})$ 的表示. 
\begin{fact}
    在共轭意义下, 存在唯一的不可约表示
    \begin{equation*}
        i: \mathrm{PSL}(2,\mathbb{R})\hookrightarrow \mathrm{PSL}(n,\mathbb{R}),
    \end{equation*}
    这个表示的构造如下: 考虑 $\mathrm{PSL}(2,\mathbb{R})$ 在 $\mathbb{R}^2 = \mR\{x,y\}$ 上的标准表示, 
    这个表示诱导了 $\mathrm{PSL}(2,\mathbb{R})$ 在 $\mathrm{Sym}^{n-1}(\mathbb{R}^2)$ 上的作用, 
    这个作用给出了一个不可约表示.
    % \begin{equation*}
    %     i: \mathrm{PSL}(2,\mathbb{R})\hookrightarrow \mathrm{PSL}(n,\mathbb{R}).
    % \end{equation*}
\end{fact}
这个嵌入将曲面基本群的每个 $\mathrm{PSL}(2,\mathbb{R})$ 表示诱导为 $\mathrm{PSL}(n,\mathbb{R})$ 表示,
\begin{equation*}
    i\circ \rho: \pi_1(S)\to \mathrm{PSL}(2,\mathbb{R})\hookrightarrow \mathrm{PSL}(n,\mathbb{R}).
\end{equation*}
由此定义了映射
\begin{equation*}
    I:\mathrm{Teich}(S)\cong \mathrm{Hom}(\pi_1(S), \mathrm{PSL}(2,\mathbb{R}))\sslash\sim\;
    \longrightarrow \mathrm{Hom}(\pi_1(S), \mathrm{PSL}(n,\mathbb{R}))\sslash\sim.
\end{equation*}
\begin{definition}
    特征簇 $\mathrm{Hom}(\pi_1(S), \mathrm{PSL}(n,\mathbb{R}))\sslash\sim$ 可能不止一个连通分支, 
    我们称包含 $I(\mathrm{Teich}(S))$ 的连通分支为 \emph{Hitchin 分支}.
\end{definition}
\begin{theorem}(Hitchin, 1992)
    \begin{equation*}
        \mathrm{Hit}_n(S) \cong \bigoplus_{i=2}^n H^0(\Sigma, K^i)\cong \mathbb{R}^{(n^2-1)(2g-2)}.
    \end{equation*}
    特别地, 
    \begin{itemize}[leftmargin=*]
        \item 当 $n=2$ 时, Hitchin 分支就是 Teichm\"uller 空间, $\mathrm{Hit}_2(S)\cong\mR^{6g-6}$; 
        \item 当 $n=3$ 时, Hitchin 分支是 $\mathrm{Hit}_3(S)\cong\mR^{16g-16}$.
    \end{itemize}
\end{theorem}

\subsection{为什么考虑 $\mathrm{PSL}(2,\mathbb{R})\hookrightarrow \mathrm{PSL}(n,\mathbb{R})$?}

\subsubsection{$\mR P^2$-结构}
\subsubsection{仿射几何, 仿射球面}
\subsubsection{完整的故事}